%% kapa: exam=1 round=25 year=2014 period=1 form=A subject=civil_law subject_ko=민법 seq=20 number=20
%% source: Q-Net 감정평가사 기출문제 (cst00309) artlSeq=5101207 file=extracted/1008/5101207/2014년 제25회 감정평가사 1차시험 기출문제/1교시(A형).hwp
%% conversion: kordoc 4.24.51 -> tools/gen1.py (auto); review=proofread-sonnet; flags=none
\begin{question}{20}
대리에 관한 설명으로 옳은 것은? (다툼이 있으면 판례에 의함)
\begin{choices}
  \item 부동산 입찰절차에서 동일한 물건에 관하여 이해관계가 다른 2인 이상의 대리인이 된 경우, 그 대리인이 한 입찰은 유효하다.
  \item 본인을 위하여 금전소비대차계약을 체결할 권한을 수여받은 대리인에게는, 특별한 사정이 없는 한, 본래의 계약을 해제할 대리권도 있다.
  \item 대리인이 수인인 경우에 다른 정함이 없으면 대리인은 공동으로 대리권을 행사하여야 한다.
  \item 피성년후견인이 아닌 자가 임의대리인이 된 후에 그에게 성년후견이 개시되면 그 대리권은 소멸한다.
  \item 대리권의 범위가 정해지지 않은 대리인은 보존행위만을 할 수 있다.
\end{choices}
\end{question}
\answer{4}
